\section{Canonical elementary principalization with compact support}

Canonical principalization can retain its marked-ideal centre rule while
exporting all the elementary geometric conclusions. Compactness of the
ambient manifold is unnecessary when the ideal's zero locus is compact.
The compatibility established here is with open restrictions and
biholomorphisms. General smooth morphisms require additional invariant
and blowup base-change theorems.

\begin{theorem}[Canonical marked resolution with compact cosupport]
\label{hir:canonical-resolution-compact-cosupport}
\uses{hir:canonical-resolution,hir:principalization-compact-support}
Let $\mathfrak m$ lie in the domain of a resolution invariant with
well-ordered values and have compact cosupport on a complex manifold.
There exists a resolution which is a canonical segment with nonempty
centres. The invariant is upper semicontinuous along canonical segments,
every open restriction is canonical, and canonical resolutions after a
biholomorphism are isomorphic, stage by stage, over that biholomorphism.
\lean{Hironaka.HironakaCanonicalPrincipalization.canonicalResolution_of_isCompact}
\leanok
\end{theorem}
\begin{proof}
Along a canonical segment the transformed cosupport is compact: it is
closed and contained in the inverse image of the initial compact cosupport
under a proper composition of blowups. Upper semicontinuity gives a
maximum on every nonempty transformed cosupport. Its maximum locus is
a nonempty smooth admissible centre with empty interior. The invariant
strictly decreases above that centre and stays unchanged elsewhere.
The compact-cosupport, well-founded driver therefore produces a
resolving segment using only nonempty maximum loci.

Upper semicontinuity is the invariant's second property. Its open
restriction property identifies every nonempty restricted centre with
the restricted maximum locus; centres missing an open set restrict to
empty centres. For two canonical resolving segments related by a
biholomorphism, induction gives corresponding maximum loci and the
uniqueness of each blowup gives isomorphisms of their stages. Their
lengths are equal: if one were longer, its prefix corresponding to the
completed shorter resolution would have empty cosupport, and a nonempty
next maximum locus could not exist. The resulting stage isomorphisms
therefore extend through both complete towers.
\end{proof}

\begin{theorem}[Canonical principalization relative to an SNC boundary]
\label{hir:canonical-elementary-principalization-boundary}
\uses{hir:canonical-resolution-compact-cosupport,hir:principalization-boundary}
Let $\mathfrak a$ be a coherent ideal on a complex manifold, nonzero
at every point, with compact zero locus, and let $E_0$ be an SNC boundary.
There is an admissible principalization $f$ obtained from the canonical
resolution of $(X,X,E_0,\mathfrak a,1)$. Its effective divisor $D$
satisfies $f^*\mathfrak a=\mathcal O(-D)$, and its accumulated boundary
$E$ is SNC. The support of $D$ is both $f^{-1}Z(\mathfrak a)$ and the
union of the new members of $E$; $f$ is an isomorphism off
$Z(\mathfrak a)$. The same tower has the open-restriction and
biholomorphism compatibility of the preceding theorem.
\lean{Hironaka.HironakaCanonicalPrincipalization.exists_canonical_principalization_boundary}
\leanok
\end{theorem}
\begin{proof}
The boundary is transverse to the whole manifold, so the marked ideal
$(X,X,E_0,\mathfrak a,1)$ lies in the invariant's domain. Its cosupport
is exactly $Z(\mathfrak a)$: order at least one is precisely vanishing
of every member of the stalk. Apply the preceding theorem to the
canonical invariant. Every centre of the resulting test tower has
normal crossings with the accumulated boundary and lies over the
initial cosupport, proving admissibility over $Z(\mathfrak a)$.

The controlled-transform identity at mark one identifies the total
pullback ideal, at the terminal empty cosupport, with the tower's
effective exceptional divisor ideal. The tower transform preserves
SNC. Induction on its blowups shows that the support of its exceptional
divisor is the union of the new boundary members. Pullback of coherent
ideals identifies that support with $f^{-1}Z(\mathfrak a)$.
The exceptional locus is contained in this support, and a proper
composition of blowups is an isomorphism over the complement of its
exceptional image. This yields the asserted isomorphism off
$Z(\mathfrak a)$. All compatibility clauses belong to this very tower,
since it was selected by the canonical marked-ideal producer.
\end{proof}
